\chapter{Calcul algébrique et polynômes}
\section*{Exercice 2.1 -- Développer et contrôler}
Développer $(2x-3)^2-(x+1)(x-1)$, puis vérifier le résultat en $x=0$ et $x=2$.

\section*{Exercice 2.2 -- Domaine avant simplification}
Simplifier $(x^2-9)/(x-3)$. Résoudre ensuite $(x^2-9)/(x-3)=6$.

\section*{Exercice 2.3 -- Racine et facteur}
Factoriser $P(x)=x^3-2x^2-x+2$ en utilisant $P(1)=0$, puis résoudre $P(x)=0$.

\section*{Exercice 2.4 -- Une identité générale}
Pour $n\ge1$, démontrer $(x-a)\sum_{k=0}^{n-1}x^{n-1-k}a^k=x^n-a^n$.

\section*{Exercice 2.5 -- Identifier des polynômes}
Deux polynômes de degré au plus $2$ coïncident en trois réels distincts. Prouver qu'ils sont égaux.

\section*{Exercice 2.6 -- Lettres de rôles différents}
Déterminer les réels $a,b$ pour lesquels $(x-a)(x-b)=x^2-5x+6$ pour tout réel $x$.
